\documentclass[10pt,a4paper]{amsart}
\usepackage[margin=19mm]{geometry}
\usepackage{amsmath,amssymb,mathtools}
\usepackage[scr=rsfs,cal=euler]{mathalfa}
\usepackage{xfrac}
\usepackage[unicode,hidelinks,bookmarks=false]{hyperref}
\newcommand{\ie}{i.e.\ }
\newcommand{\cf}{cf.\ }
\newcommand{\resp}{respectively\ }
\newcommand{\Subsection}{\subsection}
\providecommand{\nobreakdash}{\nobreak\hbox{-}}
\setlength{\emergencystretch}{2em}
\multlinegap0pt
\usepackage{xcolor}
\usepackage[shortlabels]{enumitem}
\usepackage{tikz}
\usepackage{amssymb}
\usepackage{mathtools}
\usepackage{float}
\usepackage{cancel}
\usepackage{dsfont}
\newtheorem{proposition}{Proposition}
\title{Experimenting with Permutation Wordle}
\author{Aurora Hiveley}
\date{Source: arXiv:2506.23452v2 (CC BY 4.0)}
\begin{document}
\maketitle

\noindent\emph{Source and licence.} Experimenting with Permutation Wordle. Version arXiv:2506.23452v2 (CC BY 4.0), distributed by arXiv under the Creative Commons Attribution 4.0 International licence, http://creativecommons.org/licenses/by/4.0/, as recorded in the arXiv licence metadata for this exact version and shown on the abstract page https://arxiv.org/abs/2506.23452v2. Full licence notice: \texttt{licenses/arxiv-2506.23452-CC-BY-4.0.txt}.\par\smallskip

\noindent\textit{Editorial scope.} Counted unit: the article's proposition prop:derange (source0.tex lines 147--149). The article's complete original proof of it is delivered with it (source0.tex lines 151--174, one \texttt{proof} environment of 24 lines). No mathematics is written, repaired or re-derived: the statement and the proof are the article's own lines, in the article's own notation, and no retained line is substituted or re-flowed.  No further source-local context is needed: every cross-reference of every retained line resolves inside the two delivered spans.  Nothing was omitted from any retained span, so the package carries no omission record.  No retained line contains a citation, so the wrapper carries no bibliography and no bibliography entry.  No retained line contains a figure environment, an image-inclusion command or drawing code, so no image is delivered and the package declares no asset dependency.  Editorial formatting only, in addition: an AMS \texttt{amsart} preamble that restates the article's own declarations and macros used by the retained lines, namely the environments \texttt{proposition} on separate counters, together with the standard packages those lines need.  The retained environments print in this document as Proposition 1 in order of appearance, whereas the article prints the same items with the numbering of its own full document.  The complete native source of the article as distributed in the arXiv e-print for this exact version is bundled unchanged as \texttt{sources/180708/source0.tex}.
\begin{proposition} \label{prop:derange}
    For any deranged strategy component $S[n]$, the average size of $\mathcal{J}_2$ across all games of permutation wordle whose secret permutation is a derangement of length $n$ is $\frac{n}{n-1}$.
\end{proposition}

\begin{proof}
To study a game using $S[n]$, first note that if $\gamma_1 = [1,2,\dots,n]$, then $\gamma_2 = (S[n])^{-1}$, assuming that the secret permutation is a derangement. Then $\mathcal{J}_2 = \{i \in [n] \mid d[i] = (S[n])^{-1}[i] \} $, so the set of correct positions is precisely the set of positions where the derangement and the inverse of our strategy component agree. Since the inverse of a derangement is also a derangement, and inverses are in one-to-one correspondence, we proceed using a derangement $\delta = (S[n])^{-1}$. Note then that $|\mathcal{J}_2| = \left| \{ i \in [n] \mid \delta(i) = d(i) \} \right|$ for a secret permutation $d$, so the average size of $|\mathcal{J}_2|$ across all $d \in D_n$ is $\frac{1}{|D_n|} \sum_{d \in D_n} \left| \{ i \in [n] \mid \delta(i) = d(i) \} \right|$. To prove the claim that this average is equal to $\frac{n}{n-1}$, we will show the following equivalent statement:

\begin{equation} \label{eq:derange}
    \sum_{d \in D_n} \left| \{ i \mid d(i) = \delta (i) \} \right| = \frac{n}{n-1} \cdot | D_n|
\end{equation}

\vspace{0.25cm}

On the left hand side of Eq. \ref{eq:derange}, we may equivalently fix a position in our fixed derangement $\delta$ and count the number of other derangements which share that position, then sum across all possible positions. Then we have that 
$ \sum_{d \in D_n} \left| \{ i \mid d(i) = \delta(i) \} \right| 
= \sum_{i=1}^n \left| \{ d \in D_n \mid d(i) = \delta(i) \} \right| $. From there, we may calculate $| \{ d \in D_n \mid d(i) = \delta(i) \} |$ by utilizing the recurrence for $|D_n|$. Recall that $|D_n| = (n-1)(|D_{n-1}| + |D_{n-2}|)$, wherein we construct a derangement of length $n$ by choosing a position for the element $n$ and then deranging the remaining elements. There are $n-1$ possible locations for $n$. Say that $d[n] = j$. Then if $d[j] = n$ as well, the number of ways to derange the remaining elements is $|D_{n-2}|$. If $d[j] \neq n$, then the number of ways to derange the elements of $[n] \setminus \{n\}$ is $|D_{n-1}|$ since while $j$'s fixed position is occupied, we prohibit $d[j] = n$.

In our case for each fixed entry $i$, we know its location in $\delta$ so we do not need to \textit{choose} one of $n-1$ possible locations for the element $i$. Then it suffices to calculate the number of ways to derange the remaining elements, which is $|D_{n-1}| + |D_{n-2}|$. Thus, we have the following:

\begin{align*}
    \sum_{d \in D_n} \left| \{ i \mid d(i) = \delta(i) \} \right|
    &= \sum_{i=1}^n \left| \{ d \in D_n \mid d(i) = \delta(i) \} \right| \\
    &= \sum_{i=1}^n (|D_{n-1}| + |D_{n-2}|) \\ 
    &= n \cdot (|D_{n-1}| + |D_{n-2}|) \\ 
    &= n \cdot \left( \frac{|D_n|}{n-1} \right) 
\end{align*}
    
\end{proof}

\end{document}
